\section{Proof of Theorem~\ref{thm:first}}\label{sec:local}

We keep $a=\as$. The first step is a local limit theorem for sums of independent Bernoulli variables. It is
standard; we include the short argument because it is used for $S_n$ here and for two other families of sums in
Section~\ref{sec:second}.

\begin{lemma}\label{lem:local}
For $m\ge1$ let $X_m$ be a sum of finitely many independent Bernoulli variables, with mean $\mu^{(m)}$ and variance
$V^{(m)}$, and let $d_m$ be an integer. If $V^{(m)}\to\infty$ and $d_m-\mu^{(m)}=o\bigl(\sqrt{V^{(m)}}\bigr)$, then
$\sqrt{V^{(m)}}\,\Prob(X_m=d_m)\to1/\sqrt{2\pi}$.
\end{lemma}

\begin{proof}
Let $B$ be a Bernoulli variable with parameter $q$ and $v=q(1-q)$. The centred variable $B-q$ has variance $v$ and
absolute third moment at most $v$, so for $|t|\le1$,
\[
\E e^{it(B-q)}=1-\tfrac12vt^2+O(v|t|^3),\qquad
\log\E e^{it(B-q)}=-\tfrac12vt^2+O\bigl(v|t|^3\bigr),
\]
with absolute constants. Summing over the summands of $X_m$, we get
$\log\E\exp(it(X_m-\mu^{(m)}))=-\frac12V^{(m)}t^2+O(V^{(m)}|t|^3)$ for $|t|\le1$, and hence
$\E\exp(iu(X_m-\mu^{(m)})/\sqrt{V^{(m)}})\to e^{-u^2/2}$ for every real $u$. By Fourier inversion and the
substitution $t=u/\sqrt{V^{(m)}}$,
\[
\sqrt{V^{(m)}}\,\Prob(X_m=d_m)=\frac1{2\pi}\int_{-\pi\sqrt{V^{(m)}}}^{\pi\sqrt{V^{(m)}}}
\E e^{iu(X_m-\mu^{(m)})/\sqrt{V^{(m)}}}\;e^{-iu(d_m-\mu^{(m)})/\sqrt{V^{(m)}}}\,du .
\]
The integrand, set equal to zero outside the interval of integration, converges pointwise to $e^{-u^2/2}$ and, by~\eqref{eq:cf} and $\sin(|t|/2)\ge|t|/\pi$, its modulus is
at most $e^{-2u^2/\pi^2}$. Dominated convergence gives the limit $(2\pi)^{-1}\int e^{-u^2/2}du=1/\sqrt{2\pi}$.
\end{proof}

\begin{proof}[Proof of Theorem~\ref{thm:first}]
Write $\ell=\log n$. By~\eqref{eq:muV}, $\mu_n=\as\ell+O(1)$ and $V_n=\as\ell+O(1)$, so for each fixed integer
$r\ge0$ the point $k_n+r$ satisfies the hypothesis of Lemma~\ref{lem:local} and
\[
\sqrt{\ell}\;\Prob(S_n=k_n+r)\to\frac1{\sqrt{2\pi\as}} .
\]
The left side is bounded uniformly in $r$ and $n$ by~\eqref{eq:atom}, and $\sum_r\as^{-r}<\infty$, so by dominated
convergence
\[
\sum_{r\ge0}\as^{-r}\,\Prob(S_n=k_n+r)=\frac{1+o(1)}{(1-1/\as)\sqrt{2\pi\as\,\ell}} .
\]
By~\eqref{eq:F} and Stirling's formula, $F_n(\as)=(1+o(1))\,n^{\as}/\Gamma(1+\as)$. Lemma~\ref{lem:first}
and~\eqref{eq:tilt} now give
\[
\E N_{n,k_n}=\frac n2F_n(\as)\,\as^{-k_n}\sum_{r\ge0}\as^{-r}\,\Prob(S_n=k_n+r)
=\frac{(1+o(1))\;n^{1+\as}\,\as^{-k_n}}{2\,\Gamma(1+\as)(1-1/\as)\sqrt{2\pi\as\,\ell}} ,
\]
and $n^{1+\as}\as^{-b_n}/\sqrt{\ell}=1$, as in the proof of Theorem~\ref{thm:upper}.
\end{proof}

\begin{remark}
Theorem~\ref{thm:first} is a statement about sums of Stirling numbers of the first kind in the range where the
second index is close to $\as$ times the logarithm of the first. Asymptotic expansions of these numbers go back to
Moser and Wyman~\cite{MW}; the expansion of Hwang~\cite{Hwang} is uniform for a second index up to any fixed multiple
of the logarithm of the first, and Theorem~\ref{thm:first} could also be derived from it.
\end{remark}

\begin{remark}
The convergence in Theorem~\ref{thm:first} is slow. For $a=\as$ one has
$\mu_n=\as\log n-\as\psi(1+\as)+o(1)$, where $\psi$ is the digamma function and $\as\psi(1+\as)\approx5.07$, so at
$k=b_n+O(1)$ the Gaussian factor $\exp(-(k-\mu_n)^2/(2V_n))$ of the local limit theorem differs from one by a term
of order $(\log\log n)^2/\log n$, and for moderate $n$ the constant $\as\psi(1+\as)$ is the larger part of
$k-\mu_n$. At $k=\lceil b_n\rceil$ the ratio of $\E N_{n,k}$, computed exactly
from~\eqref{eq:F}, to $C_*\as^{-(k-b_n)}$ is $0.388$, $0.458$, $0.614$ and $0.718$ for $n=10^2,10^3,10^4,10^5$; the
Gaussian factor at the same points is $0.320$, $0.436$, $0.606$ and $0.714$.
\end{remark}
